\chapter{Kết luận}
\label{sec:p1-ket-luan}
Nhóm đã tích hợp nguyên lý, D-algorithm, PODEM, full scan/trải khung,
lõi Python và mô phỏng lỗi độc lập. Trên c17, PODEM đạt $34/34$ lỗi gốc
hoặc $22/22$ đại diện; nén từ 34 xuống 7 hoặc 22 xuống 6 pattern,
cả hai tập nén phủ $34/34$ lỗi gốc. Full scan đạt $18/18$ lỗi stem;
CLI tuần tự kết hợp PODEM và vét cạn chuỗi bổ sung đạt $1/18$, $13/18$,
$18/18$ tại $k=1,2,3$ từ trạng thái đầu chưa biết.
Trace tay và code thống nhất; bốn lần chọn cube D-algorithm và hai phép
gán PI PODEM khác đơn vị, không suy ra tăng tốc.
Kết quả chỉ áp dụng stuck-at trên các mạch nhỏ, không bảo đảm mọi khuyết
tật vật lý hay hiệu năng trên mạch lớn. Hướng tiếp theo là sinh chuỗi
không phụ thuộc $Q@0$, heuristic FAN, scan compression và lỗi trễ.
